\section{Dashboard Gráfico de Calidad}
\label{sec:dashboard}

El siguiente dashboard refleja de manera visual las distribuciones derivadas de las estadísticas de ejecución consolidadas de las métricas de calidad de software implementadas. Las capturas a continuación representan los paneles configurados en Grafana para el monitoreo en tiempo real del estado de calidad del proyecto.

\subsection{Tasa de Éxito de Pruebas (TEP)}
La Tasa de Éxito de Pruebas (TEP) mide el porcentaje de casos de prueba ejecutados satisfactoriamente. Como se observa en la Figura~\ref{fig:dashboard-tep}, el sistema mantiene un 100\% de éxito en la ejecución de su suite de pruebas automatizadas (43 pruebas exitosas de 43 totales), lo que valida la robustez de las reglas de negocio implementadas en el backend.

\begin{figure}[H]
    \centering
    \includegraphics[width=0.85\textwidth]{../imagenes/tep.png}
    \caption{Panel de Monitoreo de Tasa de Éxito de Pruebas (TEP) en Grafana}
    \label{fig:dashboard-tep}
\end{figure}

\subsection{Cobertura Funcional (CF)}
La Cobertura Funcional (CF) representa la proporción de historias de usuario (HUs) clave validadas por tests automáticos. Actualmente, se tiene un 58.33\% de cobertura (7 HUs cubiertas de 12 definidas en total). Como ilustra la Figura~\ref{fig:dashboard-cf}, el panel cuenta con una tabla de desglose que identifica con precisión qué requerimientos están cubiertos y cuáles están pendientes de pruebas.

\begin{figure}[H]
    \centering
    \includegraphics[width=0.85\textwidth]{../imagenes/cf.png}
    \caption{Monitoreo de Cobertura Funcional (CF) y Desglose de Historias de Usuario}
    \label{fig:dashboard-cf}
\end{figure}

\subsection{Densidad de Defectos (DD)}
Esta métrica asocia la cantidad de errores reportados mediante commits de corrección en Git con las líneas de código del proyecto (KLOC). La Figura~\ref{fig:dashboard-dd} muestra una densidad de 1.23 defectos por cada 1,000 líneas de código, un nivel sumamente controlado respecto al estándar de la industria (menor a 2.0).

\begin{figure}[H]
    \centering
    \includegraphics[width=0.85\textwidth]{../imagenes/dd.png}
    \caption{Histórico de Densidad de Defectos (DD) y Tamaño de Código (KLOC)}
    \label{fig:dashboard-dd}
\end{figure}

\subsection{Vulnerabilidades Críticas (OWASP) (VCO)}
El panel de seguridad (Figura~\ref{fig:dashboard-vco}) consolida la auditoría de dependencias (composer audit) y el análisis estático SAST en el frontend. Se reportan 184 incidencias críticas/altas, explicadas principalmente por la inyección de código mediante el uso de \texttt{innerHTML} en el frontend sin sanitización previa y por dependencias secundarias en el backend.

\begin{figure}[H]
    \centering
    \includegraphics[width=0.85\textwidth]{../imagenes/vco.png}
    \caption{Resumen de Vulnerabilidades de Seguridad (VCO) y Distribución por Severidad}
    \label{fig:dashboard-vco}
\end{figure}

\subsection{Tiempo de Respuesta Promedio (TRP)}
El Tiempo de Respuesta Promedio (TRP) evalúa la latencia del API. La Figura~\ref{fig:dashboard-trp} muestra la evolución temporal del tiempo de respuesta agrupado por horas en las últimas 24 horas, junto con una tabla que detalla el Top 5 de los endpoints más lentos para enfocar labores de optimización y sintonía fina.

\begin{figure}[H]
    \centering
    \includegraphics[width=0.85\textwidth]{../imagenes/trp.png}
    \caption{Panel de Tiempos de Respuesta (TRP) y Rendimiento de Endpoints}
    \label{fig:dashboard-trp}
\end{figure}
